\begin{paramdefinition}[Component, determinant line and frame metrics]\label{hyp:flat-determinant-frames}
Fix a component \(Z\) of the deep cover over its coefficient field \(F\),
and a label \(\tau\) of the CM field \(E\). Write \(V_\tau,V_{\bar\tau}\)
for the two Lie eigenspaces, of dimensions \(p,q\). For frames \(v,w\) put
\[
 \mathcal L_\tau=\det V_\tau\otimes(\det V_{\bar\tau})^{-1},\qquad
 e(v,w)=(v_1\wedge\cdots\wedge v_p)\otimes
                      (w_1\wedge\cdots\wedge w_q)^{-1}.
\]
The flat weight is \((-1,+1)\) on the two determinants, with zero
theta exponent. Indeed, if \(\phi=f(v,w)e(v,w)\), then
\[
 f(vA,wB)=\frac{\det B}{\det A}f(v,w).
\]
Use the determinant metric induced by the polarization. At a finite
place take the integral eigenframes of the given prime-to-residue
level model. Write \(\mathfrak d_E\) for the different,
\(R_E=\log\operatorname{rd}(E)\), and \(D_0'=\log d_0\), where \(d_0\)
is the positive discriminant parameter of the marked quadratic field.
For the odd quadratic character \(\chi\) of \(\mathbb Q(\sqrt{-d_0})\),
use the CM normalization
\[
              \omega=\frac{L'(0,\chi)}{L(0,\chi)}+D_0'-R_E .
\]
\end{paramdefinition}

\begin{construction}
\begin{tikzpicture}[>=Stealth,every node/.style={align=center,inner sep=2pt}]
\node (line) at (0,0) {\(\mathcal L_\tau\) with generator \(\phi\)};
\node (arch) at (-2.75,-1.12) {Archimedean frames\\\(\|\phi\|_\sigma=c_\sigma\)};
\node (fin) at (2.75,-1.12) {Integral frames over \(w\)\\constant normalized order};
\node (cm) at (0,-2.35) {Compare on a CM frame};
\node (deg) at (0,-3.48) {\(\widehat{\deg}_{\rm flat}(\mathcal L_\tau)
        =\tfrac{p-q}{2}\omega+q(D'_0-R_E)\)};
\draw[->] (line) -- (arch);
\draw[->] (line) -- (fin);
\draw[->] (arch) -- (cm);
\draw[->] (fin) -- (cm);
\draw[->] (cm) -- (deg) node[midway,right=1mm] {\(\scriptstyle-\sum_w\delta_{w,\tau}\log{\rm N}w/[F:\mathbb Q]\)};
\end{tikzpicture}
\caption{Archimedean norms and finite orders of the flat generator in Definition~\ref{hyp:flat-determinant-frames}. The finite correction converts the CM frame degree to the normalized degree of Proposition~\ref{prop:deep-flat-line-degree}.}\label{fig:flat-cm-frames}
\end{construction}

\begin{proposition}[The normalized degree of the flat line]
\label{prop:deep-flat-line-degree}
For each component \(Z\) and label \(\tau\) of
Definition~\ref{hyp:flat-determinant-frames}, the determinant rule gives a
generator \(\phi\) on the proper normal generic fibre, up to scalar.
Its norm is constant on the archimedean frame tests at each embedding. At a finite place its normalized order is constant on the integral
frame tests. With the finite-order correction
\[
 \delta_{w,\tau}=q\{\operatorname{ord}_w(\tau\mathfrak d_E)
                              -\operatorname{ord}_w(d_0)\},
\]
the normalized degree is
\[
 \widehat{\deg}_{\rm flat}(\mathcal L_\tau)
       =\frac{p-q}{2}\,\omega+q(D_0'-R_E).
\]
Thus the contribution to the slope coefficient, with the squared-norm convention of this comparison, is
\[
 -\widehat{\deg}_{\rm flat}(\mathcal L_\tau)=-\tfrac{p-q}{2}\omega+q(R_E-D'_0).
\]

\end{proposition}

\begin{proof}
Choose the generic generator \(\phi\) of
Lemma~\ref{lem:flat-generic-generator}. At the embedding \(\sigma\)
write \(\Gamma,\Delta\) for the Gram matrices of the two frames.
The frame coefficient \(f\) then satisfies
\[
 |f|^2\,|\det\Gamma|/|\det\Delta|=c_\sigma^2
 \qquad(c_\sigma>0\text{ constant for the chosen }\phi).
\]
At a finite place \(w\) let \(\pi_w\) be a uniformizer. Degree zero on special-fibre curves
(Lemma~\ref{lem:flat-special-curve}) lets us apply
Lemma~\ref{lem:flat-vertical-ratio} to the projected determinant
lattice \(\det Q_\tau\). If \(R\) is the ring of a valued frame
over \(w\), the order computed in this lattice satisfies
\[
 \frac{\operatorname{ord}_{Q,R}(\phi)}{v_R(\pi_w)}
    =\max_{D\ {\rm vertical\ over}\ w}
       \frac{\operatorname{ord}_{Q,D}(\phi)}
            {\operatorname{ord}_D(\pi_w)} .
\]
To read the same section in the separate integral Lie frames,
Lemma~\ref{lem:projected-lie-frame-index} gives
\[
   \operatorname{ord}_{{\rm Lie},R}(\phi)
       =\operatorname{ord}_{Q,R}(\phi)-h_{R,\tau}.
\]
The Hodge multiplicities are fixed on this component, so
\(h_{R,\tau}/v_R(\pi_w)\) is the same at the CM frame and
at every other valued frame. The order and the metric may therefore
both be evaluated on the CM frame.
Lemma~\ref{lem:flat-cm-degree} evaluates their normalized frame
degree as \((p-q)\omega/2\).

For \(c\in F^\times\), replacing \(\phi\) by \(c\phi\) adds the
divisor of \(c\), whose normalized arithmetic degree is zero. Subtract
\(\delta_{w,\tau}\) from the order in each finite frame. Taking ideal
norms gives
\[
 \frac{1}{[F:\mathbb Q]}\sum_w\delta_{w,\tau}\log{\rm N}w
 =q\left(
 \frac{\log{\rm N}_{E/\mathbb Q}\mathfrak d_E}{[E:\mathbb Q]}
 -\log d_0\right)=q(R_E-D_0').
\]
Here the finite sum runs over the places of \(F\). The factor
\([F:E]\) from extending an ideal cancels in the normalization.
The corrected degree is the frame degree minus this sum. Its formula
contains no component-dependent quantity.
\label{cor:deep-cover-flat-degree-1}
\end{proof}

To compute the degree in Proposition~\ref{prop:deep-flat-line-degree}, we need one generator that can be evaluated in both kinds of frame. We first construct it on the generic component, then determine how its orders vary over a special fibre.

\begin{lemma}[Extension of the period generator]\label{lem:flat-generic-generator}
On the proper normal generic component of
Definition~\ref{hyp:flat-determinant-frames}, \(\mathcal L_\tau\) has an \(F\)-generator whose value in the prescribed frame at each complex embedding \(\sigma\) is a constant multiple of the normalized period determinant \(\Delta_{{\rm per},\sigma}\),
\[
 \phi_\sigma=c_\sigma\Delta_{{\rm per},\sigma},\qquad c_\sigma\ne0.
\]

\end{lemma}
\begin{proof}
In the marked period basis the first tangent block and the conjugate
of the second block have complementary directions. Pair the second
block by the Hermitian form and take the determinant. The marked
level comparison gives invariance under the change of period basis
on \(Z\). With \(\Gamma,\Delta\) the two frame Gram matrices, the
determinant identity, after division by the reference period constant,
reads
\[
 |f(v,w)|^2=\frac{|\det\Delta|}{|\det\Gamma|},\qquad
 \|f(v,w)e(v,w)\|^2=|f(v,w)|^2\frac{|\det\Gamma|}{|\det\Delta|}=1.
\]
It holds first at the embedding respecting the marked quadratic root.
Conjugating the frames preserves the absolute determinants.
\label{cor:deep-cover-flat-degree-2}

Take a proper smooth resolution of the generic component of the
normal model at \(2\), using the prime-to-\(2\) marking. In a
toroidal chart the boundary is \(z_1\cdots z_a=0\). The two
extending eigenframes are independent in the smooth semiabelian Lie
bundle. The boundary comparison for precisely these frames gives
\[
 C^{-1}Q(z)^{-A}\le u\le C Q(z)^A\quad (u=|\det\Gamma|,\ |\det\Delta|),
 \qquad Q(z)=1+\sum_{i=1}^a|\log|z_i|| .
\]
Consequently \(|f|,|1/f|\le C Q^A\), after increasing \(C,A\).
For \(g=f\) or \(1/f\), fix the other punctured coordinates \(z'\).
The Laurent coefficient of index \(-m\), \(m\ge1\), satisfies
\[
 \left|\frac{1}{2\pi i}\int_{|\zeta|=\rho}
          g(\zeta,z')\,\zeta^{m-1}\,d\zeta\right|
 \le C(z')\rho^m(1+|\log\rho|)^A\longrightarrow0 .
\]
Apply this to each boundary variable. The Laurent series on the
product of punctured discs has no negative multi-index in any
variable. Thus both \(f\) and \(1/f\) extend, and their product
remains one. Uniqueness on overlaps gives a generator on the resolution.
\label{cor:deep-cover-flat-degree-3}

Write \(Y\) for this proper geometrically integral resolution.
By GAGA and coherent base change,
\[
 H^0(Y,\mathcal L_\tau)\otimes_F\mathbb C
 \simeq H^0(Y_{\mathbb C}^{\rm an},\mathcal L_\tau^{\rm an})
 =\mathbb C\,\Phi ,
\]
where \(\Phi\) denotes the just-extended section and the line on
\(Y\) denotes its pullback. Any nonzero \(F\)-section therefore
generates after scalar extension, hence already over \(F\).
On the normal generic component use a local frame \(e\). Properness
of the resolution lifts each height-one valuation, so both
\(\phi/e\) and \(e/\phi\) have nonnegative order there. A normal
Noetherian domain is the intersection of its height-one valuation
rings. Thus both quotients are regular, as required. At another
finite prime the same construction on its prime-to-residue model
gives a generic generator. The quotient of the two generators lies
in \(F^\times\) since the generic component is proper and
geometrically integral. \label{cor:deep-cover-flat-degree-4}
\end{proof}

Fix the normal integral component of the prime-to-\(\ell\) model in
Definition~\ref{hyp:flat-determinant-frames}. Its generic point is
in the abelian locus. Let \(G\) be its smooth semiabelian group,
with its integral \(\mathcal O_E\)-action and rigidified polarizing
line. Write \(n=[E:\mathbb Q]\), \(d_G=\dim G\), and
\(r=p+q\), so \(2d_G=nr\). Over a local ring \(R\) on this
model choose formal coordinates at the identity and a rigidification
of the polarization biextension on the coordinate neighborhood.
These give the polarized tangent module
\[
       H_R=V_R\oplus V_R^*,\qquad V_R=\operatorname{Lie}(G)_R .
\]

The second summand is the space of linear terms of changes of
rigidification. Write \(U_a\) for the differential of \(a\in
\mathcal O_E\). The isomorphism between pulling the biextension
back by \(a\) in the first variable and by \(\bar a\) in the
second is represented in these frames by a unit equal to one on
both axes. Write \(C_a\) for its mixed linear coefficient, with
the first variable indexing the rows. The induced action is
\[
       (v,\alpha)\longmapsto
       (U_av,\ U_{\bar a}^{t}\alpha-C_a^{t}v).
\]

These formulas are integral even in characteristic two: the mixed
coefficient is the coefficient of \(st\) after substituting
\((sv,tw)\) with \(s^2=t^2=0\). Composition of the rigidified
isomorphisms gives matrix multiplication. A change of coordinates
acts on the first summand by its differential and on the second by
the inverse transpose, with a shear given by the mixed coefficient
of the change of frame. The inverse change has the same integral
description. Thus the displayed modules glue and commute with
base change.

\begin{auxlemma}[Local freeness over the integral coefficient algebra]
\label{lem:polarized-stalk-free}
On the marked component the polarized tangent module is locally
free of rank \(r\) over \(\mathcal O_X\otimes_{\mathbb Z}\mathcal O_E\).
In particular, on every geometric fibre of characteristic \(\ell\)
and for each Artinian local factor \(B\) of
\(k\otimes_{\mathbb Z}\mathcal O_E\),
\[
                   H_k\otimes_{k\otimes\mathcal O_E}B\simeq B^r .
\]
\end{auxlemma}
\begin{proof}
\proofstep{The abelian part of a geometric fibre.}
Work over an algebraically closed residue field \(k\) of
characteristic \(\ell\). In the interior the fibre is principally
polarized. At a boundary point choose a complete DVR dominating its stalk
whose generic point lies over the abelian generic point: normalize
the blowup of the maximal ideal, take the DVR of an exceptional
prime divisor, and complete. A one-dimensional normal stalk
itself supplies the DVR in dimension one. Its characteristic-zero
fraction field contains the local domain. After finite extension
the residue torus splits. The Raynaud extension and the cubical
line in the polarized degeneration datum reduce to
\[
  1\longrightarrow T\longrightarrow G_k
   \xrightarrow{\pi} A\longrightarrow1,\qquad
  L_{G_k}\simeq\pi^*L_A .
\]

Here the geometric fibre may be read after the further residue
field extension. Both ranks and local freeness descend faithfully.
For the period and character lattices \(Y\) and \(X^*(T)\), the
polarization induces an injection \(\Phi:Y\hookrightarrow X^*(T)\).
The degree formula for the degeneration datum gives
\[
   1=\deg\lambda_{\rm gen}
      =[X^*(T):\Phi Y]^2\deg\lambda_A,
      \qquad \deg\lambda_A=1 .
\]
The prime-to-residue model is used also when \(\ell\) divides the
original level. The Rosati adjoint is \(a\mapsto\bar a\).
\cite[Chapter~II, \S\S1--4, and Chapter~III, Corollary~7.2]{faltings-chai}

For \(0\ne a\in\mathcal O_E\) let \(K=\ker(a:A\to A)\).
The action is an isogeny when \(A\) has positive dimension, since
some integral multiple of \(a^{-1}\) is integral. The two
rigidified pullbacks of the Poincar\'e line give the perfect
Cartier pairing of \(K\) and
\(K^D=\ker(\bar a:A\to A)\).
\cite[Theorem~9.1]{milne-av-isogeny}
Let \(B_K\subset k[K]\) be the Hopf subalgebra generated by
the primitive functions, and set
\[
     J=\ker(K\longrightarrow\operatorname{Spec}B_K),\qquad
     t=\dim_k\operatorname{Lie}(K^D).
\]

A tangent vector to \(K^D\) is a character \(1+\epsilon u\) over
\(k[\epsilon]/\epsilon^2\). The character law reads
\[
             \Delta u=u\otimes1+1\otimes u .
\]
Thus the tangent vectors correspond to the primitive functions. The derivative of this function at a tangent
vector \(v\in\operatorname{Lie}K\) is the mixed coefficient of the
Cartier pairing. In our coordinates it equals
\(v^t C_a w\), up to the common choice of the inverse pairing, for
\(w\in\ker U_{\bar a}\). Indeed the two trivializations defining
the pairing on the kernels restrict to the rigidifications on the
axes. Evaluating their ratio on \((sv,tw)\) leaves just its
\(st\) coefficient.

A pair in the kernel on \(H_A\) therefore satisfies
\[
 U_av=0,\qquad U_{\bar a}^{t}\alpha=C_a^t v,\qquad
 v\in\operatorname{Lie}J .
\]
The last condition follows by pairing the middle equality with
each \(w\in\ker U_{\bar a}\). For a fixed \(v\) the ambiguity
in \(\alpha\) has dimension \(t\). For any finite group scheme
\(D\) over \(k\), its first Frobenius kernel has order
\(\ell^{\dim\operatorname{Lie}D}\), by the height-one
correspondence and the restricted enveloping algebra basis theorem.
\cite[Theorem~1.4.23]{brion-height-one}
Apply this first to \(J\), then to the Cartier dual of
\(\operatorname{Spec}B_K\), whose Lie algebra contains the \(t\)
independent primitives. Multiplicativity of orders in the finite
group scheme quotient gives
\[
   \dim\ker(a\mid H_A)\le \dim\operatorname{Lie}J+t
                          \le v_\ell(\deg a|_A).
\]
When \(A=0\) this contribution is zero.

\proofstep{The torus and the uniform kernel bound.}
The cocharacter lattice of \(T\) is torsion-free over
\(\mathcal O_E\), since multiplication by nonzero \(a\) divides
multiplication by a nonzero integer. It is projective over this
Dedekind ring. Put \(b_T=\operatorname{rk}_{\mathcal O_E}X_*(T)\)
and \(N_a=|N_{E/\mathbb Q}(a)|\).
The dimension of the kernel on \(\operatorname{Lie}T=X_*(T)\otimes k\)
is bounded by the sum of the \(\ell\)-adic elementary divisors
of multiplication by \(a\) on this lattice. For the toric term
and its conjugate dual the bounds are
\[
    \dim\ker U_{a,T}\le b_Tv_\ell(N_a),\qquad
    \dim\ker U_{\bar a,T}^t\le b_Tv_\ell(N_a) .
\]

Write \(P:V\to\operatorname{Lie}A\) for the quotient differential.
Use a rigidification pulled back from \(A\) on the geometric
fibre. The mixed coefficient there is \(P^t C_{a,A}P\).
In another rigidification the difference is a shear, so it
conjugates the action. The stable subspaces
\[
 \operatorname{Lie}T\oplus0\ \subset V\oplus P^*(\operatorname{Lie}A)^*\ \subset\ H_k
\]
have successive quotients \(\operatorname{Lie}T,H_A,
(\operatorname{Lie}T)^*\), the last with the conjugate action.
Kernel dimensions are subadditive along this filtration.
On the prime-to-\(\ell\) rational Tate module of \(A\), an
integral primitive element of \(E\) acts semisimply. Its rational
characteristic polynomial is a power of its minimal polynomial,
so every embedding occurs with multiplicity \(2\dim A/n\).
The degree of \(a\) is its determinant on this module.
\cite[Theorem~10.9 and Proposition~10.20]{milne-av-isogeny}
It follows that
\[
 \dim_k\ker(a\mid H_k)
   \le (2b_T+2\dim A/n)\,v_\ell(N_a)
   =r\,v_\ell(N_a).
\]

\proofstep{Recovering a free module at a ramified prime.}
On the characteristic-zero generic fibre the \(\tau\)-part of
\(H\) has dimension \(p+q\) for each embedding. At each stalk
the characteristic polynomial of \(a\) therefore equals that
of its regular action on \(\mathcal O_E\) raised to the \(r\)-th
power. The coefficients belong to the local domain and agree in
its fraction field, giving the integral identity
\[
      \det(X-a\mid H)=\det(X-a\mid\mathcal O_E)^{r}.
\]
For \(\mathfrak l\mid\ell\), of ramification
index \(e\) and residue degree \(f_{\mathfrak l}\), its factors
over \(k\) have the form \(B_\sigma=k[u]/(u^e)\). Each summand
\(H_\sigma\) has dimension \(er\) over \(k\). Indeed, choose
an element with nonzero residue generating the residue field at
\(\mathfrak l\) and zero residue at the other primes. The
multiplicities of its distinct nonzero eigenvalues in the preceding
characteristic polynomial are precisely these dimensions.

Choose \(a\) of valuation one at \(\mathfrak l\) and zero
at all other primes over \(\ell\), by the Chinese remainder
theorem. Then \(v_\ell(N_a)=f_{\mathfrak l}\). Multiplication
by \(a\) has the same kernel as \(u\) on the factors at
\(\mathfrak l\), and is invertible elsewhere. Thus
\[
 \sum_{\sigma}\dim_k\ker(u\mid H_\sigma)\le r f_{\mathfrak l},
 \qquad \dim_k\ker(u\mid H_\sigma)\ge \dim_k H_\sigma/e=r .
\]
All Jordan blocks have length at most \(e\), so equality forces
exactly \(r\) blocks of length \(e\) at every \(\sigma\).
This proves freeness over each \(B_\sigma\), including \(e>1\).
Over a characteristic-zero residue field the same conclusion follows
from the characteristic polynomial and separability of \(E\).

Freeness at the geometric fibre descends to projectivity at the
residue field by faithful flatness. The coefficient algebra there
is a product of local Artinian rings, and the common rank \(r\)
gives a free module. Lift a set of \(r\) coefficient generators
to the stalk. If \((b_i)\) is an integral basis of \(\mathcal O_E\),
the reductions of the \(b_i x_j\) are a basis of size \(nr=2d_G\).
Their determinant in the chosen tangent coordinates is a unit.
They give a basis over the local ring, hence the \(x_j\) are a
basis over its coefficient algebra.
\end{proof}

\begin{auxlemma}[Integral evaluation of a projected frame]
\label{lem:polarized-stalk-evaluation}
Let \(R\) be a stalk of this normal model and let the integral
character \(\sigma:\mathcal O_E\to R\) extend the label \(\tau\).
For any coefficient basis \(x_1,\ldots,x_r\) from
Lemma~\ref{lem:polarized-stalk-free}, the character quotient is
free with the indicated basis:
\[
 Q_{\sigma,R}=H_R/\sum_a(a-\sigma(a))H_R
             =\bigoplus_{j=1}^r R\,\bar x_j .
\]
Its generic-fibre realization is the image of \(H_R\) under
the characteristic-zero eigenprojection, and its bases and their
transition matrices specialize to the function field and every
valuation ring of every normalized special-fibre curve.
\end{auxlemma}
\begin{proof}
In coefficient coordinates, the map on each summand is
\(R\otimes\mathcal O_E\to R\), \(b\otimes a\mapsto b\sigma(a)\).
Its kernel is generated by \(1\otimes a-\sigma(a)\otimes1\).
After passing to the fraction field this evaluation is precisely the
corresponding projection in the split coefficient algebra. In
particular the integral quotient embeds there as the span of its
projected basis. No integral idempotent for the individual embedding
is required at a ramified prime.

Let \(C\) be a normalized curve, \(K=k(C)\) and
\(V_z=\mathcal O_{C,z}\). If \(x\) and \(x'\) are frames taken
at the image of \(z\) and of the generic point, compare them over
the latter stalk. Write their change of basis as \(T\) over its
coefficient algebra, in the convention \(x=x'T\). Then
\[
       \bar x=\bar x'\,T_\sigma,\qquad
       T_\sigma=(1\otimes\sigma)(T),\qquad
       \det T_\sigma\in\mathcal O_{X,\eta_C}^{\times}.
\]
The two maps from \(\mathcal O_{X,z}\) (at the image of \(z\))
to \(K\), through \(V_z\) and through \(\mathcal O_{X,\eta_C}\),
agree. Hence \(Q_{\sigma,V_z}\) in the generic curve frame
is \(T_\sigma V_z^r\), with the entries read in \(K\).
The determinant lattice is \((\det T_\sigma)V_z\). The determinant
can have nonzero \(z\)-valuation. Its valuation is exactly the
transition contribution to the degree on the curve.
\end{proof}

The integral image of an eigenprojection and the intersection with an
eigenspace may differ at a ramified prime. The following determinant
calculation compares the two conventions in
Proposition~\ref{prop:deep-flat-line-degree}.

\begin{auxlemma}[The local change from the projected to the two Lie frames]
\label{lem:projected-lie-frame-index}
Let \(R\) be the valuation ring of a finite extension of the local
coefficient field containing all embeddings of \(E\), write
\(K=\operatorname{Frac}R\), and let \(v_R\) be its integral valuation.
For a marked integral test write \(L=\operatorname{Lie}(G)_R\),
\(S_\gamma=L\cap(L_K)_\gamma\) and
\(n_\gamma=\dim_K(L_K)_\gamma\).
Let \(I_\tau\) be the embeddings inducing the same prime
\(\mathfrak l\mid\ell\) and residue embedding as \(\tau\).
For a uniformizer \(\theta\) of \(E_{\mathfrak l}\) put
\[
 h_{R,\tau}=\sum_{\gamma\in I_\tau\setminus\{\tau\}}
        n_\gamma v_R\bigl(\gamma(\theta)-\tau(\theta)\bigr).
\]
In the generic determinant identification, the lattice
\(\det S_\tau\otimes(\det S_{\bar\tau})^{-1}\) is contained in
\(\det Q_{\tau,R}\) with colength \(h_{R,\tau}\). In particular
the normalized difference of the two frame orders depends only on the
coefficient field, the place and the multiplicities.
\end{auxlemma}

\begin{proof}
On the factor selected by the common residue embedding the coefficient
algebra is monogenic. Indeed the unramified part splits over \(R\),
and a uniformizer generates the totally ramified integer ring over
its unramified part. On this factor put \(D=\theta-\tau(\theta)\).
It has the presentation
\[
 R[X]\big/\prod_{\gamma\in I_\tau}(X-\gamma(\theta)),\qquad
          Q_{\tau,R}=H_R/DH_R
\]
where \(H_R\) on the right denotes the selected factor. Take the
same factor of the exact sequence formed by the second summand and
the projection onto the first,
\[
         0\longrightarrow F\longrightarrow H_R
                              \longrightarrow L\longrightarrow0 .
\]
Here \(L\) in the sequence also means its selected factor, and
\(F\) is the corresponding factor of the integral dual of the
full Lie module, with conjugate action. Its projected image is
\(S_{\bar\tau}^*\): restriction to the saturated free sublattice
\(S_{\bar\tau}\) is surjective over a DVR, and the generic
restriction is the dual eigenprojection.

Since \(Q_{\tau,R}\) is free, the image of \(F/DF\) in it is its
torsion-free quotient, which the preceding description identifies with
\(S_{\bar\tau}^*\). Right exactness gives
\[
     0\longrightarrow S_{\bar\tau}^*\longrightarrow Q_{\tau,R}
                       \longrightarrow L/DL\longrightarrow0 .
\]
Write \(I=L\cap\bigoplus_{\gamma\ne\tau}(L_K)_\gamma\) on this
factor. The torsion in \(L/DL\) is \(I/DL\) and the free quotient
is \(L/I\). The colength in the assertion is consequently
\[
    \operatorname{length}(I/DL)+
    \operatorname{length}\bigl(L/(I+S_\tau)\bigr).
\]

This follows by a basis change over the DVR, first
saturating \(S_{\bar\tau}^*\) in \(Q_{\tau,R}\) and then taking
the index of \(S_\tau\) in the free quotient \(L/I\).
The map \(D\) identifies \(L/(I+S_\tau)\) with \(DL/DI\),
since its kernel on \(L\) is \(S_\tau\). The two lengths add to
\[
  \operatorname{length}(I/DI)
     =v_R\bigl(\det(D\mid I_K)\bigr)=h_{R,\tau}.
\]
The formula holds also for an unramified factor, when the sum is empty.
Passing to a further finite extension multiplies both the colength
and \(v_R(\pi_w)\) by its ramification index.
\end{proof}

\begin{lemma}[Degree on special-fibre curves]\label{lem:flat-special-curve}
Use the integral eigenframes of Definition~\ref{hyp:flat-determinant-frames}
on the normal model at a finite place of residue characteristic \(\ell\).
The determinant line measuring the order of the section \(\phi\)
of Lemma~\ref{lem:flat-generic-generator} has degree zero on the normalization \(\widetilde C\) of every complete integral curve \(C\) in this special fibre,
\[
 \deg_{\widetilde C}(\mathcal L_\tau|_{\widetilde C})=0.
\]

\end{lemma}
\begin{proof}
Apply Lemma~\ref{lem:polarized-stalk-evaluation} on the marked
model after extending constants to contain the integral label
\(\sigma\). On the generic fibre the projected polarized tangent
space is an extension of \(V_\tau\) by \(V_{\bar\tau}^*\).
Lift the \(v_i\) to this extension and adjoin the dual frame of
the \(w_j\). Its determinant is independent of the lifts and
transforms by \(\det A/\det B\). It identifies
\[
       \det Q_{\sigma,\operatorname{Frac}R}
                 =\det V_\tau\otimes(\det V_{\bar\tau})^{-1} .
\]
Its extension on the curve has the integral determinant lattices
computed in Lemma~\ref{lem:polarized-stalk-evaluation}. These
are the orders to which Lemma~\ref{lem:flat-vertical-ratio} will
be applied. The comparison with the separate Lie frames in the CM
calculation keeps the different at the given place.
\label{cor:deep-cover-flat-degree-5}

\proofstep{Compare the two Hodge degrees on a Siegel model.}
Enlarge the finite field of constants \(k\) to split the residue
fields of \(\mathcal O_E\) above \(\ell\). Write \(e\) for a primitive
idempotent in \(k\otimes\mathcal O_E\), and write \(H_e=eH\).
We compare determinants using the underlying principally polarized
family. Choose a full Siegel marking of some order \(n\ge3\) prime
to \(\ell\) on a finite cover of the abelian generic locus. Take the
normalized closure of the graph of its map to a projective Siegel
toroidal compactification at that level. If \(X\) denotes our
model and \(X'\) this graph, the projections give
\[
 X\ \longleftarrow\ X'\ \longrightarrow\ \overline{\mathcal A}_{d_G,n}^{\rm tor}.
\]
The first projection is proper and surjective. On this normal cover the two
semiabelian families agree: their generic isomorphism and its inverse
extend uniquely. \cite[I, Proposition~2.7]{faltings-chai}

There is a complete integral curve upstairs dominating
\(\widetilde C\) after normalization. Indeed a closed point of the
preimage of its generic point has residue field finite over
\(k(\widetilde C)\); take the closure in the proper fibre product
and normalize. For this finite map \(h:C'\to\widetilde C\),
\[
              \deg_{C'} h^*\det H_e
                    =(\deg h)\deg_{\widetilde C}\det H_e ,
\]
using degrees relative to the same ground field. We can therefore
compute on \(C'\), even for a curve contained in the boundary.

Let \(D\) be the pullback of the covariant logarithmic de Rham
bundle on the Siegel compactification. Its Hodge and conjugate
filtrations form an \(F\)-zip: the two filtrations are by subbundles
and their graded bundles match in reverse order after Frobenius
pullback. The boundary assertion here is the logarithmic one for
principally polarized Siegel families in characteristic \(\ell\),
including \(\ell=2\). In a toroidal chart its covariant realization
is the Lie algebra of the period-equivariant universal vector
extension of the Raynaud extension. For split torus \(T\), abelian
quotient \(A\) and period lattice \(Y\), its weight filtration has
successive quotients
\[
          \operatorname{Lie}T,\qquad D(A),\qquad Y\otimes\mathcal O .
\]
Here \(D(A)\) denotes the ordinary de Rham homology of \(A\). \cite[VI, Theorem~1.1]{faltings-chai}
\cite[\S2D, pp.~476--477]{vdg-looijenga-boundary-zip}

The coefficient action on these filtrations also extends. On the
completed normal degeneration chart, after splitting the torus,
the action preserving the polarization in the Rosati sense acts
integrally on the Raynaud extension and its period homomorphism
\(u\). If \(a_Y,a_{G^\natural}\) are the induced maps, then
\[
 u:Y\longrightarrow G^\natural_\eta,\qquad
          u\circ a_Y=a_{G^\natural}\circ u .
\]
This is the equivalence for integral endomorphism structures in
\cite[Theorem~5.1.1.4]{lan-pel-degeneration}.
Applying the universal vector extension and its Lie functor gives
the action on \(D\). It commutes with relative Frobenius and
Verschiebung by functoriality on the degeneration data. The
actions agree on overlaps by their common generic action. In
particular the construction applies to every \(a\in\mathcal O_E\),
including elements of positive \(\ell\)-adic valuation.

Put \(V=\operatorname{Lie}(G)|_{C'}\), and let \(U=V^*\) with the
conjugate dual action induced by principal polarization. Write
\(F\) for absolute Frobenius of \(C'\) and \(F_k(e)\) for the
idempotent obtained by raising the scalar coefficients of \(e\)
to the \(\ell\)-th power. Taking the indicated direct summands
of the two filtrations gives exact sequences on the whole curve:
\[
\begin{aligned}
 0&\longrightarrow U_e\longrightarrow D_e
                 \longrightarrow V_e\longrightarrow0,\\
 0&\longrightarrow F^*V_e\longrightarrow D_{F_k(e)}
                 \longrightarrow F^*U_e\longrightarrow0 .
\end{aligned}
\]
The quotients are vector bundles, so pullback to the curve and
projection by the idempotents preserve exactness.

The polarized
tangent bundle itself has \(0\to U_e\to H_e\to V_e\to0\).
Consequently degree additivity and \(\deg F^*L=\ell\deg L\)
for a line bundle give
\[
 \deg H_{F_k(e)}=\deg D_{F_k(e)}
       =\ell(\deg U_e+\deg V_e)=\ell\deg H_e .
\]
The permutation \(e\mapsto F_k(e)\) has finite order \(s>0\).
Iteration gives \((\ell^s-1)\deg H_e=0\). This argument uses
the whole local Artinian factor selected by \(e\), before reduction
by its maximal ideal.

The local Artinian factor \(B\) at the chosen character has residue
field \(k\). Filter \(B\) by ideals with simple quotients \(k\).
If \(b=\operatorname{length}_B B\), Lemma~\ref{lem:polarized-stalk-free} identifies
each quotient of the induced filtration of \(H_e\) with the same
bundle \(H_e\otimes_B k\) (up to a constant one-dimensional factor).
By degree additivity,
\[
 0=\deg H_e=b\,\deg(H_e\otimes_B k).
\]
The determinant of the eigencharacter quotient therefore has degree
zero, including when \(B\) has nilpotents.
\label{cor:deep-cover-flat-degree-6}
\end{proof}

\begin{auxlemma}[A curve through a prescribed point]
\label{lem:flat-specialization-curve}
Let \(T\) be an integral proper variety of positive dimension over
a finite field, \(U\subset T\) a dense open subset, and \(z\) a closed
point of \(T\). Let \(k\) be the finite ground field. There is a finite extension \(k'/k\), a complete integral curve \(C\subset T_{k'}\) and a point \(z'\in C\) over \(z\) such that its generic point satisfies
\[
 z'\in C\cap(T_{k'})_z,\qquad \eta_C\in U_{k'}.
\]

\end{auxlemma}
\begin{proof}
Over an algebraic closure take a reduced irreducible component through
a lift \(z'\) of \(z\). The inverse image of \(U\) is dense there.
Take an affine neighborhood \(V\) of \(z'\) on this component and a
finite Noether normalization \(f:V\to\mathbb A^d\).
For the closed subset \(Z=V\setminus U_{\bar k}\) one has
\[
   \dim f(Z)=\dim Z<d\qquad(Z\ne\varnothing).
\]
Choose a line through \(f(z')\) and a point outside \(f(Z)\).
For any component \(C_0\) of its inverse image through \(z'\),
the \(d-1\) linear equations and finiteness give
\[
 1=d-(d-1)\le \dim C_0=\dim f(C_0)\le1 .
\]
The finite image is the whole line, so \(C_0\) meets the given open
set. Its closure in the proper variety is the required integral curve.
The component, line and curve are defined by finitely many data over
the algebraic closure of a finite field and hence descend to a finite
extension.
\end{proof}

\begin{lemma}[The vertical order on a connected fibre]\label{lem:flat-vertical-ratio}
Let \(\mathcal L\) be the integral determinant line on the proper
normal component over a DVR with uniformizer \(\pi\) and connected
special fibre. Assume it has degree zero on every normalized
complete integral curve in that fibre. If \(\phi\) generates on
the generic fibre, then, for every valued frame over this DVR,
\[
 \operatorname{ord}(\phi)=\frac{a}{m}\operatorname{ord}(\pi),\qquad
 \frac{a}{m}=\max_{D\ {\rm vertical}}
            \frac{\operatorname{ord}_D(\phi)}{\operatorname{ord}_D(\pi)},
 \quad m>0 .
\]
\end{lemma}
\begin{proof}
The rational section \(s=\pi^a\phi^{-m}\) of \(\mathcal L^{-m}\)
has horizontal order zero. Its order along a vertical prime \(D\) is
\[
 a\,\operatorname{ord}_D(\pi)-m\,\operatorname{ord}_D(\phi)\ge0 .
\]
By normality \(s\) is regular. Its unit locus \(U\) contains the
generic point of a component attaining the maximum.

Take any irreducible special-fibre component \(T\) meeting \(U\).
If \(T\setminus U\) is nonempty, choose a closed point \(z\) there.
If \(\dim T=0\) then \(T\) consists of a single point, contrary to
the choice of \(z\). Otherwise apply Lemma~\ref{lem:flat-specialization-curve}
to \(T_{\rm red}\) and its open unit locus, and normalize the resulting
curve. Its image in the original fibre is an integral curve. Pullback
along the finite map of normalized curves preserves degree zero.

On this complete normal curve \(C\), the pullback of \(s\) is
regular and generically nonzero, so
\[
 0=-m\deg_C\mathcal L=\deg_C\operatorname{div}(s)
       =\sum_{x\in C_{\rm cl}}[k(x):k_C]\operatorname{ord}_x(s),
 \qquad \operatorname{ord}_x(s)\ge0 ,
\]
where \(k_C\) is its constant field. A point over \(z\) would
contribute a positive summand. Hence \(T\subset U\). Across any
intersection with a neighboring component its generic point belongs
to the same open \(U\). Connectedness of the finite component graph
gives the result on the whole fibre. Pulling the unit \(s\) back to
a valued frame gives \(m\operatorname{ord}(\phi)=a\operatorname{ord}(\pi)\).
Changing the generic generator by \(c\in F^\times\) translates its
order by \(\operatorname{ord}(c)\). The local models and framed
uniformization supply nonempty tests at the required places.
\label{cor:deep-cover-flat-degree-7}
\end{proof}

The vertical calculation has made the normalized orders independent of the component of a connected fibre. We can therefore evaluate the remaining degree at the prescribed CM frame.

\begin{lemma}[CM calculation in the frame metrics]\label{lem:flat-cm-degree}
For the polarized determinant line on \(Z\) in
Definition~\ref{hyp:flat-determinant-frames}, before the finite-order
correction of Proposition~\ref{prop:deep-flat-line-degree}, the normalized arithmetic degree in the CM frame norms satisfies
\[
 \widehat{\deg}_{\rm CM}(\mathcal L_\tau)=\frac{p-q}{2}\,\omega.
\]

\end{lemma}
\begin{proof}
At a pure CM test on the component use the two elliptic types of
\(\mathbb Q(\sqrt{-d_0})\), with their good integral abelian models,
integral Lie generators and normalized principal polarizations. Write
\(n_{\pm,s}\) for the squared Lie lengths at a complex embedding \(s\).
In a common finite field of definition \(k/F\), write \(\langle
\cdot\rangle_\infty\) for the mean over all \([k:\mathbb Q]\) complex
embeddings, counting conjugates. The Chowla--Selberg comparison and
the functional equation of the odd primitive quadratic character give
\[
 \begin{aligned}
 -\langle\log n_{\pm,s}\rangle_\infty=X
 &=\gamma+\log(2\pi/\sqrt{d_0})-L'(1,\chi)/L(1,\chi)\\
 &=L'(0,\chi)/L(0,\chi)+D_0'/2,\qquad
 \omega=X-R_E+D_0'/2 .
 \end{aligned}
\]
The mean is invariant under further field extension.
\cite{colmez-chowla-selberg}\label{cor:deep-cover-flat-degree-8}

Let \(T\) change the reference period basis to the two orthogonal
CM blocks. Let \(G\) be the reference Hermitian determinant, and
\(G_+,G_-\) the two block determinants. In these frames the
coefficient of \(\phi\) is \(\eta\,\tau(G_-)/\tau(\det T)\).
The elliptic period scalars for the two blocks are included in
\(\eta\). Evaluating the algebraic generator on the algebraic
frames shows that \(\eta\in k^\times\) after finite extension.
The identity \(G_+G_-=(\det T)\overline{\det T}\,G\) cancels the
transition determinant against the determinant metric. With the
two squared length scales \(n_{\pm,s}\sqrt{d_0}\), this gives
\[
 A_s=\log|s\eta|+\tfrac12\log|s(\tau G)|
       +\tfrac p2(\log n_{+,s}+D_0'/2)
       -\tfrac q2(\log n_{-,s}+D_0'/2).
\]
\label{cor:deep-cover-flat-degree-9}

At \(\ell\) take a locally integral CM basis with
\(t_\ell=\det T_\ell\). For \(j:k\hookrightarrow\mathbb C_\ell\)
put \(i=j\circ\tau\), \(D_j=\log|i(\mathfrak d_E)|_\ell\) and
\(d_\ell=\log|d_0|_\ell\). The absolute value of an ideal here is
that of a local generator. Unimodularity of the polarized period
lattice and trace duality with the different give
\[
 \log|i(t_\ell)|_\ell+\log|i(\bar t_\ell)|_\ell+\log|i(G)|_\ell
       =-(p+q)D_j,\qquad \log|i(G_-)|_\ell=-qD_j .
\]

The second identity uses the orthogonal pure-type block. The
locally integral period basis spans the CM product lattice over
the marked quadratic order up to prime-to-\(\ell\) index.
Projecting to each integral eigenlattice by the trace-dual basis
contributes a generator of \(\mathfrak d_E/(\sqrt{-d_0})\), whose
log absolute value is \(D_j-d_\ell/2\) on either conjugate
label. The two frame exponents are \(-p,+q\), giving the finite log norm
\[
 B_j=\log|j\eta|_\ell-\log|i(t_\ell)|_\ell-qD_j
                              +(q-p)(D_j-d_\ell/2).
\]
\label{cor:deep-cover-flat-degree-10}

To sum, put \(B'_j=B_j-q(D_j-d_\ell)\). Conjugation permutes
the embeddings of \(E\) with equal multiplicities at \(\ell\),
so the two sums for \(t_\ell,\bar t_\ell\) coincide. The first
trace-dual identity yields
\[
 \sum_{j\mid\ell} B'_j
 =\sum_{j\mid\ell}\left(\log|j\eta|_\ell+\tfrac12\log|i(G)|_\ell
              -\tfrac{p+q}{2}D_j+\tfrac{p+q}{2}d_\ell\right).
\]
Write \(\langle\cdot\rangle_f\) for the sum over finite embeddings,
divided by \([k:\mathbb Q]\), with \(|\ell|_\ell=\ell^{-1}\).
Only finitely many places contribute. Ideal norms give
\(\langle D_j\rangle_f=-R_E\), \(\langle d_\ell\rangle_f=-D_0'\).
The product formula cancels the \(\eta,G\) terms. Therefore
\[
 \begin{aligned}
 -\langle A_s\rangle_\infty-\langle B'_j\rangle_f
 &=\tfrac{p-q}{2}(X-D_0'/2)+\tfrac{p+q}{2}(D_0'-R_E)\\
 &=\tfrac{p-q}{2}\omega+q(D_0'-R_E),\\
 -\langle A_s\rangle_\infty-\langle B_j\rangle_f
 &=\tfrac{p-q}{2}\omega .
 \end{aligned}
\]
The last equality restores \(B_j-B'_j=q(D_j-d_\ell)\) in the
frame norms. \label{cor:deep-cover-flat-degree-11}
\end{proof}
